\documentclass[10pt,a4paper]{amsart}
\usepackage[margin=19mm]{geometry}
\usepackage{amsmath,amssymb,mathtools}
\usepackage[scr=rsfs,cal=euler]{mathalfa}
\usepackage{xfrac}
\usepackage[unicode,hidelinks,bookmarks=false]{hyperref}
\newcommand{\ie}{i.e.\ }
\newcommand{\cf}{cf.\ }
\newcommand{\resp}{respectively\ }
\newcommand{\Subsection}{\subsection}
\providecommand{\nobreakdash}{\nobreak\hbox{-}}
\setlength{\emergencystretch}{2em}
\multlinegap0pt
\usepackage{amssymb}
\usepackage{dsfont}
\usepackage{enumerate}
\usepackage{tikz}
\newtheorem{theorem}{Theorem}
\title{A general framework for level continuous fuzzy-valued functions}
\author{J. J. Font, S. Macario, M. Sanchis}
\date{Source: arXiv:2603.26854v1 (CC BY 4.0)}
\begin{document}
\maketitle

\noindent\emph{Source and licence.} A general framework for level continuous fuzzy-valued functions. Version arXiv:2603.26854v1 (CC BY 4.0), distributed by arXiv under the Creative Commons Attribution 4.0 International licence, http://creativecommons.org/licenses/by/4.0/, as recorded in the arXiv licence metadata for this exact version and shown on the abstract page https://arxiv.org/abs/2603.26854v1. Full licence notice: \texttt{licenses/arxiv-2603.26854-CC-BY-4.0.txt}.\par\smallskip

\noindent\textit{Editorial scope.} Counted unit: the article's theorem th:bounded1 (source0.tex lines 558--562). The article's complete original proof of it is delivered with it (source0.tex lines 563--639, one \texttt{proof} environment of 77 lines). No mathematics is written, repaired or re-derived: the statement and the proof are the article's own lines, in the article's own notation, and no retained line is substituted or re-flowed.  No further source-local context is needed: every cross-reference of every retained line resolves inside the two delivered spans.  Nothing was omitted from any retained span, so the package carries no omission record.  No retained line contains a citation, so the wrapper carries no bibliography and no bibliography entry.  No retained line contains a figure environment, an image-inclusion command or drawing code, so no image is delivered and the package declares no asset dependency.  Editorial formatting only, in addition: an AMS \texttt{amsart} preamble that restates the article's own declarations and macros used by the retained lines, namely the environments \texttt{theorem} on separate counters, together with the standard packages those lines need.  The retained environments print in this document as Theorem 1 in order of appearance, whereas the article prints the same items with the numbering of its own full document.  The complete native source of the article as distributed in the arXiv e-print for this exact version is bundled unchanged as \texttt{sources/197306/source0.tex}.
\begin{theorem}\label{th:bounded1}
Let $K$ be a compact space and
let $f\in \overline{C}([0,1]\times K)$.
Then $f$ is a bounded function.
\end{theorem}

\begin{proof}
Fix $(\lambda_0,t_0)\in [0,1]\times K$. If $\lambda_0\neq 1$ then, by the existence of right limit, we have that there exists
$\delta_{(\lambda_0,t_0)}^+ >0$ with
\begin{align}\label{eqRL1}
|f(\beta,t_0)-f(\lambda_0+,t_0)|<1 \quad \text{for $\lambda_0<\beta<\lambda_0 +\delta_{(\lambda_0,t_0)}^+$}.
\end{align}
If $\lambda_0\neq 0$, then, by the jointly left continuity - continuity at $(\lambda_0, t_0)$, we get $\delta_{(\lambda_0,t_0)}^->0$ and an open neighbourhood $V_{(\lambda_0,t_0)}$ of $t_0$
with
\begin{align}\label{eqRL2}
|f(\beta,s)-f(\lambda_0,t_0)|<1 \quad \text{for $\lambda_0-\delta_{(\lambda_0,t_0)}^-<\beta\leq \lambda_0$ and $s\in V_{(\lambda_0,t_0)}$.}
\end{align}
%Also, by right continuity - continuity at $(0,t_0)$, we get $\delta_{(0,t_0)}^+>0$ and an open neighbourhood $V_{(0,t_0)}$ of $t_0$, such that
%\begin{align}\label{eqRL22}
%|f(\beta,s)-f(0,t_0)|<1 \quad \text{for $0<\beta<\delta_{(\lambda_0,t_0)}^+$ and $s\in V_{(0,t_0)}$.}
%\end{align}

Also we can assume, by the continuity of the cross-sections $f(\lambda_0 +\delta_{(\lambda_0,t_0)}^+,\cdot)$ ($\lambda_0\neq 1$) and
$f(\lambda_0 -\delta_{(\lambda_0,t_0)}^-,\cdot)$ ($\lambda_0\neq  0$) at $t_0$, that there exists $V'_{(\lambda_0,t_0)}$ with
\begin{align}\label{eqRL3}
|f(\lambda_0 +\delta_{(\lambda_0,t_0)}^+,s)-f(\lambda_0 +\delta_{(\lambda_0,t_0)}^+,t_0)|&<1 \quad \text{for $s\in V'_{(\lambda_0,t_0)}$},\\
\label{eqRL4}
|f(\lambda_0 -\delta_{(\lambda_0,t_0)}^-,s)-f(\lambda_0 -\delta_{(\lambda_0,t_0)}^-,t_0)|&<1 \quad \text{for $s\in V'_{(\lambda_0,t_0)}$}.
\end{align}
Take $W_{(\lambda_0,t_0)}=V_{(\lambda_0,t_0)}\cap V'_{(\lambda_0,t_0)}$ and  define the following open sets in $[0,1]\times K$:
\begin{align*}
U_{(\lambda_0,t_0)}&:=(\lambda_0 -\delta_{(\lambda_0,t_0)}^-,\lambda_0 +\delta_{(\lambda_0,t_0)}^+)\times W_{(\lambda_0,t_0)}, \lambda_0\neq 0,1;\\
U_{(0,t_0)}&:=[0,\lambda_0 +\delta_{(0,t_0)}^+)\times W_{(0,t_0)};\\
U_{(1,t_0)}&:=(1 -\delta_{(\lambda_0,t_0)}^-,1]\times W_{(1,t_0)}.
\end{align*}

 Fix $(\beta,s)\in U_{(\lambda_0,t_0)}$. We have two possibilities
depending on which side of $\lambda_0$ is $\beta$.

\medskip
If $\lambda_0-\delta_{(\lambda_0,t_0)}^-<\beta\leq \lambda_0$ and $s\in W_{(\lambda_0,t_0)}$, by \eqref{eqRL2}, we have
\begin{align*}
|f(\beta,s)| &\leq |f(\beta,s)-f(\lambda_0,t_0)|+|f(\lambda,t_0)|<1
+|f(\lambda_0,t_0)|=\rho_{(\lambda_0,t_0)}^-.
\end{align*}

If  $\lambda_0<\beta<\lambda_0 +\delta_{(\lambda_0,t_0)}^+$ and $s\in W_{(\lambda_0,t_0)}$,  we  need the monotonicity for controlling the behaviour of $f(\beta,s)$.

Firstly,  we assume that $f(\cdot,s)$ is nondecreasing.
By \eqref{eqRL3} and \eqref{eqRL1}, we have
\begin{align*}
f(\beta,s) &\leq f(\lambda_0 +\delta_{(\lambda_0,t_0)}^+,s)=f(\lambda_0 +\delta_{(\lambda_0,t_0)}^+,s)-f(\lambda_0 +\delta_{(\lambda_0,t_0)}^+,t_0)\\
&+ f(\lambda_0 +\delta_{(\lambda_0,t_0)}^+,t_0 )-f(\lambda_0+,t_0)+f(\lambda_0+,t_0)<1+1+f(\lambda_0+,t_0)=q_{(\lambda_0,t_0)}^+.
\end{align*}
and, by \eqref{eqRL2},
\begin{align*}
f(\beta,s) &\geq f(\lambda_0,s)=f(\lambda_0,s)-f(\lambda_0,t_0)+ f(\lambda_0,t_0 )>-1+f(\lambda_0,t_0)=q_{(\lambda_0,t_0)}^-.
\end{align*}

Then,
$$q_{(\lambda_0,t_0)}^-\leq f(\beta,s)\leq q_{(\lambda_0,t_0)}^+$$

If $f(\cdot,s)$ is nonincreasing, then by reversing the inequalities, we have
\begin{align*}
f(\beta,s) &\geq -1-1+f(\lambda_0+,t_0)=p_{(\lambda_0,t_0)}^-,\\
f(\beta,s) &\leq 1+f(\lambda_0,t_0)=p_{(\lambda_0,t_0)}^+.
\end{align*}
and we obtain the bounds
$$p_{(\lambda_0,t_0)}^-\leq f(\beta,s)\leq p_{(\lambda_0,t_0)}^+.$$

Taking $\rho_{(\lambda_0,t_0)}^+=\max\{|q_{(\lambda_0,t_0)}^-|, |q_{(\lambda_0,t_0)}^+|,|p_{(\lambda_0,t_0)}^-|,|p_{(\lambda_0,t_0)}^+|\}$
we obtain, for every $(\beta,s) \in U_{(\lambda_0,t_0)}$,
$$|f(\beta,s)|\leq\rho_{(\lambda_0,t_0)}:=\max\{\rho_{(\lambda_0,t_0)}^-,\rho_{(\lambda_0,t_0)}^+\}.$$

Now,
 $\{  U_{(\lambda_0,t_0)}\}_{(\lambda_0,t_0)\in[0,1]\times K}$ is an open cover of the compact space $[0,1]\times K$. Then, we can find a finite subcover, say
 $\{  U_{(\lambda_k,t_k)}\}_{k=1}^{N}$.
 Let $\rho=\max\{ \rho_{(\lambda_k,t_k)}\ : 1\leq k\leq N\}$. Hence,
 for every $(\beta,s)\in [0,1]\times K$, we can find $1\leq k\leq N$ such that
 $(\beta,s)\in U_{(\lambda_k,t_k)}$, which means
 $$|f(\beta,s)|\leq\rho_{(\lambda_k,t_k)}\leq \rho$$
 and we are done.
\end{proof}

\end{document}
